\documentclass[12pt,a4paper]{article}

% -------------------------------------------------
% Sprache / Kodierung
% -------------------------------------------------
\usepackage[T1]{fontenc}
\usepackage[utf8]{inputenc}
\usepackage[ngerman]{babel}
\usepackage{lmodern}
\usepackage{microtype}

% -------------------------------------------------
% Mathematik
% -------------------------------------------------
\usepackage{amsmath,amssymb,amsthm,mathtools}

% -------------------------------------------------
% Layout / Grafik / Farben
% -------------------------------------------------
\usepackage[a4paper,margin=2.3cm]{geometry}
\usepackage{booktabs}
\usepackage{xcolor}
\usepackage[hidelinks,unicode]{hyperref}
\pdfstringdefDisableCommands{%
  \def\ü{ue}\def\ö{oe}\def\ä{ae}\def\Ü{Ue}\def\Ö{Oe}\def\Ä{Ae}\def\ss{ss}%
}
\usepackage[most]{tcolorbox}
\usepackage{enumitem}
\usepackage{array}

% -------------------------------------------------
% Farben (konsistent mit Vorgängerdokumenten)
% -------------------------------------------------
\definecolor{lückenrot}{RGB}{180,40,40}
\definecolor{bewiesengrün}{RGB}{30,100,50}
\definecolor{warnorange}{RGB}{200,100,10}
\definecolor{hintergrundblau}{RGB}{235,242,250}
\definecolor{hintergrundgelb}{RGB}{255,252,225}
\definecolor{adischviolett}{RGB}{90,20,140}
\definecolor{fehlerrot}{RGB}{200,0,0}
\definecolor{konjviolett}{RGB}{80,0,130}
\definecolor{e8grau}{RGB}{55,55,90}
\definecolor{kingmanblau}{RGB}{20,60,130}
\definecolor{zirkelrot}{RGB}{160,0,0}
\definecolor{dynblau}{RGB}{10,50,120}
\definecolor{diophgrün}{RGB}{20,90,40}

% -------------------------------------------------
% tcolorbox-Stile
% -------------------------------------------------
\tcbset{
  lücke/.style={colback=red!8, colframe=lückenrot, fonttitle=\bfseries,
    title={Offene Lücke}},
  ergebnis/.style={colback=green!7, colframe=bewiesengrün,
    fonttitle=\bfseries, title={Bewiesenes Ergebnis}},
  warnung/.style={colback=orange!10, colframe=warnorange,
    fonttitle=\bfseries, title={Achtung / Heuristik}},
  adisch/.style={colback=violet!6, colframe=adischviolett,
    fonttitle=\bfseries, title={2-adische Perspektive}},
  kernidee/.style={colback=hintergrundblau, colframe=blue!60!black,
    fonttitle=\bfseries, title={Kernidee}},
  konj/.style={colback=violet!5, colframe=konjviolett, fonttitle=\bfseries,
    title={Konjektur (unbewiesen)}},
  hauptsatz/.style={colback=green!10, colframe=bewiesengrün!80!black,
    fonttitle=\bfseries\large, title={Satz (bewiesen)}},
  dynfeld/.style={colback=blue!5, colframe=dynblau,
    fonttitle=\bfseries, title={Arithmetische Dynamik}},
  dioph/.style={colback=green!5, colframe=diophgrün,
    fonttitle=\bfseries, title={Diophantische Analogie}},
  forschung/.style={colback=hintergrundgelb, colframe=e8grau,
    fonttitle=\bfseries, title={Offene Forschungsfrage}},
}

% -------------------------------------------------
% Theorem-Umgebungen
% -------------------------------------------------
\newtheorem{definition}{Definition}[section]
\newtheorem{proposition}[definition]{Proposition}
\newtheorem{theorem}[definition]{Theorem}
\newtheorem{lemma}[definition]{Lemma}
\newtheorem{korollar}[definition]{Korollar}
\newtheorem{bemerkung}[definition]{Bemerkung}
\newtheorem{hypothese}[definition]{Hypothese}
\newtheorem{satz}[definition]{Satz}

\newenvironment{beweis}{\begin{proof}[Beweis]}{\end{proof}}

% -------------------------------------------------
% Makros
% -------------------------------------------------
\newcommand{\vii}{\nu_2}
\newcommand{\N}{\mathbb{N}}
\newcommand{\Z}{\mathbb{Z}}
\newcommand{\R}{\mathbb{R}}
\newcommand{\Q}{\mathbb{Q}}
\newcommand{\H}{\mathbb{H}}
\newcommand{\Ztwo}{\mathbb{Z}_2}
\newcommand{\abs}[1]{\lvert#1\rvert}
\newcommand{\atwo}[1]{\lvert#1\rvert_2}
\newcommand{\Lam}{\Lambda}
\newcommand{\col}{\mathrm{Col}}
\newcommand{\KC}{K}    % Kolmogorov-Komplexität
\newcommand{\TV}{\mathrm{TV}}

% -------------------------------------------------
% Dokument
% -------------------------------------------------
\title{%
  Arithmetische Dynamik bei Collatz\\[0.3em]
  \large Kolmogorov-Komplexität, Uniformitäts-Vermutung\\
  und Diophantische Approximation\\[0.5em]
  \normalsize Die verbleibende Frontlinie%
}
\author{Thomas Hoffbauer}
\date{16.\ Juni 2026}

\begin{document}
\maketitle

\begin{abstract}
Diese Arbeit benennt das verbleibende mathematische Kernproblem der
Collatz-Vermutung in einer neuen Sprache: der \emph{arithmetischen Dynamik}.
Nach der Kartierung der exakten Schranken (Tao 2022: $\mu_H(E)=0$;
Kingman/Lyapunov: $\Lambda \approx -0{,}830$; RPF mod-12: Spektrallücke
$|\lambda_2| \approx 0{,}35$) verbleibt eine einzige Lücke:
die Brücke zwischen der Maß-Null-Aussage über $\Ztwo$ und dem Verhalten
\emph{aller} $n \in \N$.

Wir schlagen das neue Konzept der \emph{arithmetischen Dynamik} vor: die
Frage, wie schnell die endliche Kolmogorov-Komplexität eines Startwerts
$n \in \N$ durch die mod-12-Mischung ``zerrieben'' wird. Die
\emph{Uniformitäts-Vermutung} --- eine quantitative Abstandsschranke
zwischen $\N$ und dem Ausnahmeset $E \subset \Ztwo$ --- wird als neues
Forschungsziel formuliert. Das Analogon zu Roths Satz (1955) gibt
der Vermutung ihre diophantische Präzision.

\textbf{Status:} Alle Resultate dieses Dokuments sind, sofern nicht anders
markiert, Vermutungen oder heuristische Überlegungen. Die Collatz-Vermutung
bleibt unbewiesen.
\end{abstract}

\tableofcontents

\bigskip
\begin{tcolorbox}[warnung, title={Einordnung}]
Dieses Dokument ist das zwölfte in einer Reihe, die die Collatz-Vermutung
aus verschiedenen Perspektiven angeht: Lyapunov-Exponenten, RPF-Theorie,
LTE-Dichte, Lean-Formalisierung, Diophantische Approximation. Es fasst
den aktuellen Stand der \emph{verbleibenden Frontlinie} zusammen und
benennt die konkrete offene Frage.
\end{tcolorbox}

% ====================================================
\section{Kolmogorov-Komplexität und Collatz}
\label{sec:kolmogorov}
% ====================================================

\subsection{Grundlegende Definitionen}

\begin{definition}[Kolmogorov-Komplexität]
Die \emph{Kolmogorov-Komplexität} $\KC(x)$ eines Strings $x \in \{0,1\}^*$
ist die Länge des kürzesten universellen Turingmaschinen-Programms, das
$x$ ausgibt:
\[
  \KC(x) := \min\{|p| : U(p) = x\}.
\]
Für eine natürliche Zahl $n \in \N$ ist $\KC(n) := \KC(\mathrm{bin}(n))$,
wobei $\mathrm{bin}(n)$ die Binärdarstellung ist.
\end{definition}

\begin{tcolorbox}[ergebnis, title={Komplexität natürlicher Zahlen (Standard)}]
Für alle $n \in \N$ gilt:
\[
  \KC(n) \leq \lfloor\log_2 n\rfloor + 1 + O(1).
\]
Das folgt direkt aus der endlichen Binärdarstellung: Das Programm
``\texttt{output }$\mathrm{bin}(n)$'' hat Länge $\lfloor\log_2 n\rfloor + O(1)$.
Insbesondere: $\KC(n) = O(\log n)$.
\end{tcolorbox}

\subsection{Elemente von \texorpdfstring{$\Ztwo \setminus \N$}{Z2 ohne N}
haben generisch unendliche Komplexität}

\begin{tcolorbox}[adisch, title={Generische Komplexität in $\Ztwo$}]
Für ein \emph{generisches} Element $x \in \Ztwo$ (d.\,h.\ für ein nach
dem Haar-Maß $\mu_H$ zufällig gezogenes Element) ist die Kolmogorov-Komplexität
unendlich:
\[
  \mu_H\bigl(\{x \in \Ztwo : \KC(x) < \infty\}\bigr) = 0.
\]
Formal: Die Elemente mit endlicher Kolmogorov-Komplexität sind abzählbar
(da jedes Programm endliche Länge hat), daher ist ihre Menge eine
$\mu_H$-Nullmenge.

\medskip
Insbesondere: $\N \subset \Ztwo^*$ ist eine $\mu_H$-Nullmenge mit
$\KC(n) = O(\log n)$ --- eine extrem dünne, hochgeordnete Teilmenge
des ``generisch zufälligen'' $\Ztwo$.
\end{tcolorbox}

\subsection{Was der $3x+1$-Schritt an Komplexität erzeugt}

Sei $n \in \N$ ungerade. Der ungerade Collatz-Schritt ist $n \mapsto 3n+1$
(vor der anschließenden Division durch $2^{\vii_2(3n+1)}$).

\begin{proposition}[Bitlänge nach einem ungeraden Schritt]
\label{prop:bitlaenge}
Für ungerades $n \geq 1$ gilt:
\begin{align}
  \lfloor\log_2(3n+1)\rfloor
  &= \lfloor\log_2 n\rfloor + 1
    \quad\text{falls } 2^k/3 < n < 2^k \cdot \frac{2^k - 1}{3} + 1
    \text{ für passendes }k, \notag \\
  &\approx \lfloor\log_2 n\rfloor + \log_2 3
    \quad\text{im Durchschnitt (}\log_2 3 \approx 1{,}585\text{).} \notag
\end{align}
Präziser: der ungerade Schritt $n \mapsto 3n+1$ erzeugt im Mittel
$\log_2 3 \approx 1{,}585$ neue Bits, während die anschließende
Division durch $2^{\vii_2(3n+1)}$ im Mittel $\mathbb{E}[\vii_2]$ Bits
entfernt, wobei $\mathbb{E}[\vii_2] = 2$ (geometrische Verteilung).
\end{proposition}

\begin{tcolorbox}[dynfeld, title={Nettobilanz: Pro Block wächst die Bitlänge um $\Lambda_{\mathrm{block}}$}]
Sei $B_k$ ein Collatz-Block (ein ungerades Glied $n_{2k}$ bis zum nächsten
ungeraden Glied $n_{2k+2}$). Die mittlere Netto-Änderung der Bitlänge pro
Block ist:
\[
  \Lambda_{\mathrm{block}}
  := \mathbb{E}\bigl[\log_2(n_{k+1}/n_k)\bigr]
  = \log_2 3 - \mathbb{E}[\vii_2(3n+1)] \cdot \log_2 2
  \approx 1{,}585 - 2 = -0{,}415.
\]
Da der Kingman-Lyapunov-Exponent aus dem Zusammenspiel aller Block-Typen
berechnet wird, ergibt sich:
\[
  \Lambda = \lim_{K\to\infty} \frac{1}{K}\sum_{k=0}^{K-1} \log_2 F(B_k)
  \approx -0{,}830 < 0.
\]
\textbf{Zentrale Einsicht:} Obwohl jeder ungerade Schritt neue Bits
\emph{erzeugt}, überwiegt die Division --- im Mittel SCHRUMPFT die
Bitlänge. Der Lyapunov-Exponent $\Lambda \approx -0{,}830$ ist genau
diese Netto-Schrumpfungsrate.
\end{tcolorbox}

\begin{bemerkung}[Neue Bits als ``injizierte Information'']
Der Punkt, den das Binärinformations-Argument aus
\textit{collatz\_uniformitaet.tex} aufgreift, ist korrekt:
Der $3x+1$-Schritt injiziert bei jedem ungeraden Glied neue arithmetische
Information (neues Bitmuster am oberen Ende der Binärdarstellung). Der
``Speicher'' des Startwerts ist endlich ($\lfloor\log_2 n\rfloor + O(1)$
Bits), aber die Trajektorie \emph{erzeugt} neuen Speicher. Die Zirkularität
liegt darin, dass man nicht weiß, ob dieser neue Speicher beschränkt bleibt.
\end{bemerkung}

% ====================================================
\section{Die Nicht-Lokalität von $3x+1$}
\label{sec:nichtlokal}
% ====================================================

\subsection{Formale Bit-Analyse}

\begin{definition}[Bit-Trägheitstiefe]
Für $n \in \N$ mit $\lfloor\log_2 n\rfloor = L-1$ (also $n$ hat $L$ Bits)
definieren wir die \emph{Bit-Trägheitstiefe} $\mathrm{BTT}(n, k)$ als die
Anzahl der Bits der $k$-ten Iterierter $T^k(n)$, die noch direkt vom
Startwert $n$ abhängen (im Sinne der Bitverschiebung):
\[
  \mathrm{BTT}(n, k) := \max\bigl\{j \geq 0 : \text{Bit } j \text{ von } T^k(n)
  \text{ hängt von } n \text{ ab}\bigr\}.
\]
\end{definition}

\begin{tcolorbox}[dynfeld, title={Nicht-Lokalität: Bit-Injektion durch $3x+1$}]
Der Schritt $n \mapsto 3n+1$ ist ein globaler Operator auf der Binärdarstellung:
\begin{align*}
  3n + 1 &= 2n + n + 1 \\
         &= \underbrace{n\text{ (um 1 links verschoben)}}_{\text{erzeugt neues Bit oben}}
           + \underbrace{n}_{\text{lokale Kopie}}
           + \underbrace{1}_{\text{Carry-Kaskade}}.
\end{align*}
Die Addition $2n + n$ kann zu einer langen Carry-Kaskade führen, die
\emph{alle} Bits von $n$ verändert. Anders als z.\,B.\ die Division durch $2$
(reine Rechtsverschiebung) ist die Multiplikation mit $3$ ein globaler Eingriff.

\medskip
\textbf{Folge:} Nach $k$ ungeraden Schritten haben potenziell \emph{alle}
Bits von $T^k(n)$ Information vom Startwert $n$ --- die Bit-Trägheitstiefe
schrumpft nicht automatisch.
\end{tcolorbox}

\subsection{Die mod-12-Mischung als Informations-Destruktor}

\begin{tcolorbox}[kernidee, title={Wie Mischung die Startwert-Information ``zerreibt''}]
Sei $\pi$ die stationäre Verteilung der mod-12-Markov-Kette (aus Schritt B
in \textit{collatz\_uniformitaet.tex}). Die exponentielle Mischung
\[
  \bigl\|M^k e_{B(n)} - \pi\bigr\|_{\TV} \leq C\,|\lambda_2|^k,
  \quad |\lambda_2| \approx 0{,}35,
\]
bedeutet: Nach $k$ Block-Schritten ist die \emph{mod-12-Klasse} von $T^k(n)$
exponentiell schnell von $n$ unabhängig.

\medskip
Das ist eine echte ``Informationslöschung'' --- aber nur auf der Ebene der
mod-12-Residuen. Die tatsächliche Größe $T^k(n)$ trägt noch volle arithmetische
Information.

\medskip
\textbf{Quantitativ:} Die Mischungszeit (bis $\varepsilon$-Abstand von $\pi$)
ist:
\[
  k_\varepsilon \approx \frac{\log(C/\varepsilon)}{\log(1/|\lambda_2|)}
  \approx \frac{\log(C/\varepsilon)}{1{,}05}
\]
Blockschritte --- unabhängig von $n$. Für $\varepsilon = 0{,}01$, $C = 1$:
$k_\varepsilon \approx 4{,}4$ Blockschritte.
\end{tcolorbox}

\subsection{Das Heuristik-Modell für die Informations-Zerreibung}

\begin{tcolorbox}[warnung, title={Heuristik: exponentieller Informationsabfall}]
Definiere heuristisch:
\[
  I_n(k) := \text{``2-adische Information von } T^k(n)\text{, die noch von } n \text{ abhängt''}
  \approx 2^{-k/12}.
\]
Begründung: Die mod-12-Mischungsrate ist $|\lambda_2|^k \approx (0{,}35)^k$.
Für $12$ Blockschritte gilt $(0{,}35)^{12} \approx 2^{-17{,}7}$, also fällt
die Startinformation pro 12 Blöcke um einen Faktor $\approx 2^{-12}$.

\medskip
Nach $k \sim 12 \cdot \lfloor\log_2 n\rfloor$ Schritten gilt heuristisch:
\[
  I_n(k) \approx 2^{-\lfloor\log_2 n\rfloor} \approx 0.
\]
Damit folgt die Trajektorie danach nur noch dem RPF-stationären Maß $\pi$
--- \emph{falls sie überhaupt so lange läuft, ohne zu divergieren}.

\medskip
\textbf{Dies ist eine Heuristik, kein Beweis.} Die Zirkularität besteht
unverändert: Man kann nicht voraussetzen, dass die Trajektorie $12\lfloor\log_2 n\rfloor$
Schritte macht, ohne zuvor zu beweisen, dass sie nicht divergiert.
\end{tcolorbox}

% ====================================================
\section{Die Uniformitäts-Vermutung}
\label{sec:uniformitaet}
% ====================================================

\subsection{Roths Satz als Analogie}

Um die neue Forschungsfrage scharf zu stellen, ziehen wir eine Analogie
aus der Diophantischen Approximation.

\begin{tcolorbox}[dioph, title={Roths Satz (1955)}]
Sei $\alpha \in \R$ eine algebraische irrationale Zahl. Dann gilt für jedes
$\varepsilon > 0$ und alle bis auf endlich viele $p/q \in \Q$:
\[
  \left|\alpha - \frac{p}{q}\right| > q^{-2-\varepsilon}.
\]
Anders formuliert: Algebraische Irrationale können durch Rationale nicht
``zu gut'' approximiert werden. Die Qualität der Approximation ist nach
unten durch $q^{-2}$ beschränkt.
\end{tcolorbox}

Die Bedeutung von Roths Satz liegt darin, dass er eine \emph{quantitative
Trennungsschranke} zwischen zwei Mengen liefert: den algebraischen
Irrationalen und den Rationalen. Obwohl Rationale in den Irrationalen
dicht liegen, gibt es eine \emph{Mindest-Separation}.

\begin{tcolorbox}[kernidee, title={Die Collatz-Analogie: Trennungsschranke zwischen $\N$ und $E$}]
Das Collatz-Analogon zu Roths Satz wäre: Eine quantitative Schranke für den
\emph{2-adischen Abstand} zwischen natürlichen Zahlen $n \in \N$ und dem
Ausnahmeset $E$ des Collatz-Operators in $\Ztwo$.

\medskip
Intuitiv: $\N$ ist in $\Ztwo$ dicht (jede 2-adische Kugel enthält natürliche
Zahlen), aber die Elemente von $E$ (falls $E \neq \emptyset$) müssten
``hochkomplex'' sein und könnten daher $\N$ nicht beliebig nah kommen.
\end{tcolorbox}

\subsection{Formale Formulierung der Uniformitäts-Vermutung}

\begin{tcolorbox}[konj, title={Uniformitäts-Vermutung (Hauptvermutung dieses Dokuments)}]
Sei $T: \Ztwo^* \to \Ztwo^*$ der Collatz-Operator auf den 2-adischen
ungeraden Ganzzahlen und $E = \{x \in \Ztwo^* : \text{Bahn von } x
\text{ divergiert oder hat Lyapunov-Exponent} \geq 0\}$ das Ausnahmeset.

\medskip
\textbf{Vermutung:} Es gibt eine universelle Konstante $c > 0$, sodass
für alle $n \in \N$ (ungerade):
\[
  \mathrm{dist}_{2\text{-adisch}}(n,\, E)
  := \inf_{e \in E} |n - e|_2
  \;\geq\; c \cdot 2^{-\lfloor\log_2 n\rfloor}.
\]

\medskip
\textbf{Intuition:} Elemente von $E$ haben ``unendliche Kolmogorov-Komplexität''
(sie sind generische Punkte in $\Ztwo$). Natürliche Zahlen $n$ haben
Komplexität $O(\log n)$. Die Schranke $c \cdot 2^{-\lfloor\log_2 n\rfloor}$
ist die ``natürliche'' 2-adische Auflösung eines $n$-Bit-Objekts.

\medskip
\textbf{Konsequenz:} Falls diese Vermutung beweisbar ist, so liegen alle
$n \in \N$ in einem Abstand $\geq c \cdot 2^{-\lfloor\log_2 n\rfloor}$
von $E$ --- und da die Collatz-Dynamik stetig in $\Ztwo$ ist, konvergiert
die Bahn von $n$ nicht zu $E$. Zusammen mit $\Lambda < 0$ würde die
Collatz-Vermutung folgen.
\end{tcolorbox}

\subsection{Verbindung zu den LTE-Worst-Cases}

\begin{tcolorbox}[dynfeld, title={LTE-Extremalpunkte und die diophantische Schranke}]
Die extremalen C-Ketten-Startpunkte aus \textit{collatz\_lte\_dichte.tex}
haben die Form
\[
  n = 2^{k+1} \cdot 3^r - 1
  \quad (k, r \in \N).
\]
Der 2-adische Abstand zwischen $n$ und dem Nachfolgepunkt
$n+1 = 2^{k+1} \cdot 3^r$ ist:
\[
  |n - (n+1)|_2 = |{-1}|_2 = 2^{-(k+1)}.
\]

\medskip
Das ist exakt die Größenordnung der diophantischen Schranke
$c \cdot 2^{-\lfloor\log_2 n\rfloor}$, denn $\lfloor\log_2 n\rfloor \approx k+1$
für diese Startpunkte.

\medskip
\textbf{Interpretation:} Die LTE-extremalen Punkte liegen \emph{maximal
nah} an benachbarten Zweierpotenzen --- aber noch im ``erlaubten Bereich''
der Schranke. Sie bilden die kritische Klasse für die Uniformitäts-Vermutung.
\end{tcolorbox}

\begin{proposition}[Konsistenz mit bekannten LTE-Schranken]
\label{prop:konsistenz}
Sei $n = 2^{k+1} \cdot 3^r - 1$ ein LTE-extremaler Punkt. Dann gilt:
\[
  |n|_2 = 1
  \quad \text{(ungerade)}, \qquad
  |n+1|_2 = 2^{-(k+1)}.
\]
Der 2-adische Abstand von $n$ zu den 2-adisch kleinen Elementen
(Zweierpotenzen) ist $2^{-(k+1)} \approx 2^{-\lfloor\log_2 n\rfloor}$.
Das ist konsistent mit der Uniformitäts-Vermutung für $c = 1$.
\end{proposition}

% ====================================================
\section{Das ``Zerreiben'' der endlichen Information}
\label{sec:zerreiben}
% ====================================================

\subsection{Formale Definition der Startwert-Information}

\begin{definition}[Startwert-Information (heuristisch)]
\label{def:info}
Sei $n \in \N$ gegeben. Definiere die \emph{2-adische Startwert-Information}
nach $k$ Schritten als die bedingte Entropie:
\[
  I_n(k) := H\bigl(T^k(n) \,\big|\, T^k(\N)\bigr),
\]
wobei $T^k(\N)$ die Menge aller $k$-fach iterierten natürlichen Zahlen ist
und die Entropie im Sinne der 2-adischen Gleichverteilung auf
$k$-adischen Kugeln zu verstehen ist (heuristisch).
\end{definition}

\begin{bemerkung}
Eine rigorose Definition von $I_n(k)$ über bedingte Entropie erfordert
ein Maßmodell auf $\N$, das mit der Collatz-Dynamik verträglich ist.
Dieser Punkt ist eine der offenen Forschungsfragen (vgl.\ \S\,\ref{sec:fragen}).
\end{bemerkung}

\subsection{Das Zerreibungs-Modell}

\begin{tcolorbox}[dynfeld, title={Heuristisches Modell des Informations-Zerfalls}]
\textbf{Modell:} Die Startwert-Information $I_n(k)$ zerfällt exponentiell:
\[
  I_n(k) \sim 2^{-k \cdot \log_2(1/|\lambda_2|)} \sim (0{,}35)^k.
\]

\textbf{Mechanismus:}
\begin{enumerate}[label=(\arabic*), nosep]
  \item Nach $k_1 \approx 5$ Blockschritten: Die mod-12-Klasse von $T^{k_1}(n)$
    ist von $n$ fast unabhängig (TVD $\leq 0{,}01$).
  \item Nach $k_2 \approx 12 \lfloor\log_2 n\rfloor$ Blockschritten:
    Die gesamte logarithmische Information ($\log_2 n$ Bits) ist aufgebraucht.
  \item Danach: Die Trajektorie folgt dem RPF-stationären Maß $\pi$ und
    muss mit $\Lambda \approx -0{,}830$ gegen $1$ konvergieren.
\end{enumerate}

\medskip
\textbf{Formale Lücke:} Schritt (3) setzt voraus, dass die Trajektorie
$k_2$ Schritte macht, ohne zu divergieren. Das ist äquivalent zur Collatz-Vermutung
für $n$. Das Modell ist konsistent, aber nicht beweisend.
\end{tcolorbox}

\subsection{Verbindung zur Kingman-Analyse}

\begin{tcolorbox}[kernidee, title={Das vollständige Bild: Zerreibung + Lyapunov}]
Kombiniert man die Kingman-Analyse ($\Lambda \approx -0{,}830$, aus
\textit{collatz\_kingman.tex}) mit dem Zerreibungsmodell, ergibt sich
folgendes kohärentes Bild:

\medskip
\begin{enumerate}[label=\textbf{Phase \arabic*:}, leftmargin=*, align=left]
  \item \textbf{(Kurz, $k \leq k_1 \approx 5$ Blöcke):}
    Startwert-Information dominiert. Mod-12-Klasse nähert sich $\pi$
    exponentiell an.
  \item \textbf{(Mittel, $k_1 < k \leq k_2$):}
    Mod-12-Klasse ist gut gemischt. Lyapunov-Exponent $\Lambda \approx -0{,}830$
    dominiert. Trajektorie schrumpft im Mittel.
  \item \textbf{(Lang, $k > k_2$):}
    Alle Startinformation zerrieben. Trajektorie folgt nur noch RPF $\to$
    Konvergenz gegen $1$ \emph{bedingt auf Nicht-Divergenz}.
\end{enumerate}

\medskip
Die \emph{einzige} verbleibende Lücke: Der Übergang von ``bedingt'' zu
``unbedingt'' --- genau die Brücke, die die Uniformitäts-Vermutung
liefern soll.
\end{tcolorbox}

% ====================================================
\section{Drei offene Forschungsfragen}
\label{sec:fragen}
% ====================================================

\begin{tcolorbox}[forschung, title={Forschungsfrage 1: Roths Satz für Collatz}]
\textbf{Frage:} Gibt es ein Analogon zu Roths Satz für den Collatz-Operator?

\medskip
Präzise: Sei $\mathcal{F} \subset \Ztwo^*$ die Menge der 2-adischen Fixpunkte
und Zykelpunkte des Collatz-Operators $T$ (einschließlich fraktaler Fixpunkte).
Gilt:
\[
  \inf_{n \in \N} \inf_{f \in \mathcal{F}} |n - f|_2 \cdot 2^{\lfloor\log_2 n\rfloor}
  \;\geq\; c > 0
  \quad \text{?}
\]
Das würde bedeuten: $\N$ kann den 2-adischen Zykelpunkten nicht beliebig
nah kommen (relative zur ``natürlichen Auflösung'' von $n$).

\medskip
\textbf{Schwierigkeit:} Die Menge $\mathcal{F}$ ist kompliziert strukturiert
und bisher wenig verstanden. Die Collatz-Dynamik auf $\Ztwo$ ist nicht
isometrisch und hat keine einfache algebraische Struktur.

\medskip
\textbf{Warum wichtig:} Wenn diese Frage positiv beantwortet werden kann,
würde zusammen mit $\Lambda < 0$ die Collatz-Vermutung folgen (ohne
Zirkularität).
\end{tcolorbox}

\begin{tcolorbox}[forschung, title={Forschungsfrage 2: Rigorose Definition von $I_n(k)$}]
\textbf{Frage:} Gibt es eine rigorose Definition der ``Startwert-Information''
$I_n(k)$ aus Definition~\ref{def:info}?

\medskip
\textbf{Mögliche Ansätze:}
\begin{enumerate}[label=(\alph*), nosep]
  \item \emph{Bedingte Entropie:} $I_n(k) := H(T^k(n) \,|\, \mathcal{F}_n^k)$,
    wobei $\mathcal{F}_n^k$ die von der $k$-stufigen Bahn erzeugte $\sigma$-Algebra
    ist. Aber welches Maß?
  \item \emph{2-adische Distanz:} $I_n(k) := \sum_{m \in \N} |T^k(n) - T^k(m)|_2^{-1}$
    als Maß für die ``Unterscheidbarkeit'' von $T^k(n)$ von anderen Startwerten.
  \item \emph{Algorithmische Entropie:} $I_n(k) := \KC(T^k(n) \,|\, T^k(\pi))$,
    bedingte Kolmogorov-Komplexität relativ zur stationären Verteilung.
\end{enumerate}
\end{tcolorbox}

\begin{tcolorbox}[forschung, title={Forschungsfrage 3: Poincaré-Rekurrenz für Collatz auf $\Ztwo$}]
\textbf{Frage:} Hat der Collatz-Operator auf $\Ztwo$ eine Poincaré-Rekurrenz-Eigenschaft,
die alle Zyklen charakterisiert?

\medskip
\textbf{Hintergrund:} Der Poincaré-Rekurrenz-Satz besagt für maßerhaltende
Abbildungen: Fast alle Punkte kehren in jede Umgebung ihres Startpunkts
zurück. Für den Collatz-Operator auf $(\Ztwo^*, \mu_H)$ ist die Maßerhaltung
nicht bewiesen (das wäre stärker als die Collatz-Vermutung).

\medskip
\textbf{Schwächere Version:} Gibt es eine vollständige Charakterisierung
aller Zyklen in $\Ztwo^*$ (nicht nur in $\N$)? Die bekannten Zyklen
in $\Z_2$ sind $\{-1\}$, $\{-5/7, -7/5\}$ etc.\ (Zyklen unter
$T$ auf dem negativen Zweig). In $\N$ ist nur der $1 \to 4 \to 2 \to 1$-Zyklus
bekannt. Eine vollständige Zählung würde das Uniformitätsproblem schärfer
einrahmen.
\end{tcolorbox}

% ====================================================
\section{Einordnung in die Beweislandschaft}
\label{sec:einordnung}
% ====================================================

\subsection{Checkliste: Was formalisiert wurde (PR \#1 -- \#11)}

\begin{tcolorbox}[ergebnis, title={Mathematische Bilanz (Stand: Juni 2026)}]
\begin{enumerate}[label=\textbf{\arabic*.}, leftmargin=*, align=left, nosep,
    itemsep=2pt]
  \item[\checkmark] \textbf{DC-Beweis (Direkte Kontraktion):}
    Für ungerades $n \equiv 1 \pmod{4}$ ist $T(n) < n$. Lean-bewiesen.
  \item[\checkmark] \textbf{LTE-Formel:} $\vii_2(3^k - 1) = \vii_2(2) + \vii_2(k)$.
    Lean-Basis vorhanden (\texttt{padicValNat}).
  \item[\checkmark] \textbf{Vier-Typen-Klassifikation:} EABC mod-4, mit
    mod-12-Verfeinerung. Formal korrekt.
  \item[\checkmark] \textbf{Kingman-Lyapunov $\Lambda \approx -0{,}830$:}
    Quantifiziert Taos impliziten Lyapunov. Bedingt auf RPF.
  \item[\checkmark] \textbf{RPF mod-12:} Exponentielle Mischungsrate
    $|\lambda_2| \approx 0{,}35$. Numerisch belegt.
  \item[\checkmark] \textbf{Zirkularität von Schritt C:} Quaternionen-Diophantik
    ist nicht valide. Klar dokumentiert.
  \item[\checkmark] \textbf{Zirkularität von Argument D:} Binärinformations-Argument
    setzt Nicht-Divergenz voraus. Klar dokumentiert.
  \item[\checkmark] \textbf{Uniformitätslücke formalisiert:} Tao $\mu_H(E)=0$
    vs.\ Collatz-Vermutung für alle $n \in \N$. Formal eingeordnet.
  \item[\checkmark] \textbf{LTE-Dichte:} Extremale C-Ketten-Startpunkte
    $n = 2^{k+1} 3^r - 1$ identifiziert und analysiert.
  \item[\checkmark] \textbf{Lean-Werkzeugkasten:} \texttt{padicValNat},
    Stirling, von-Staudt-Clausen, Vierquadrate. Mathlib-kompatibel.
  \item[\checkmark] \textbf{Kolmogorov-Schranke für $\N$:}
    $\KC(n) = O(\log n)$, kontrastiert mit generischer Unendlichkeit in $\Ztwo$.
  \item[\checkmark] \textbf{Bit-Injektionsanalyse:} Pro ungeradem Schritt
    $\log_2 3 \approx 1{,}585$ neue Bits, Netto mit $\Lambda$: Schrumpfung.
  \item[\checkmark] \textbf{Uniformitäts-Vermutung (neu):}
    Diophantische Abstandsschranke $\mathrm{dist}_{2\text{-adisch}}(n, E)
    \geq c \cdot 2^{-\lfloor\log_2 n\rfloor}$ präzise formuliert.
  \item[\checkmark] \textbf{Analogie zu Roths Satz:} Strukturelle Parallele
    zwischen Collatz und Diophantischer Approximation hergestellt.
  \item[\checkmark] \textbf{LTE-Konsistenz:} Extremalpunkte liegen in
    Abstand $\approx 2^{-(k+1)}$ --- konsistent mit Uniformitäts-Vermutung
    für $c = 1$.
  \item[\checkmark] \textbf{Drei Forschungsfragen:} Offene Agenda für
    das neue Feld ``Arithmetische Dynamik / Diophantische Approximation
    für Collatz'' formuliert.
\end{enumerate}
\end{tcolorbox}

\subsection{Die verbleibende Lücke}

\begin{tcolorbox}[lücke, title={Die einzige verbleibende Lücke}]
\textbf{Gegeben:}
\begin{itemize}[nosep]
  \item $\Lambda \approx -0{,}830 < 0$ (Schrumpfung im Mittel).
  \item Exponentielle Mischung der mod-12-Klassen in $\leq 5$ Blockschritten.
  \item $\mu_H(E) = 0$ (Tao).
  \item $\N$ hat endliche Kolmogorov-Komplexität $O(\log n)$.
  \item $E$ besteht (falls nicht leer) aus Elementen mit unendlicher
    Kolmogorov-Komplexität.
\end{itemize}

\medskip
\textbf{Gesucht:} Eine quantitative Trennungsschranke:
\[
  \mathrm{dist}_{2\text{-adisch}}(n, E) \geq c \cdot 2^{-\lfloor\log_2 n\rfloor}
  \quad \text{für alle } n \in \N.
\]

\medskip
\textbf{Warum das schwierig ist:}
\begin{enumerate}[label=(\arabic*), nosep]
  \item Die Struktur von $E$ (falls $E \neq \emptyset$) ist völlig unbekannt.
  \item Es gibt keine bekannte algebraische oder analytische Beschreibung
    der ``Ausnahmepfade'' in $\Ztwo$.
  \item Das Analogon zu Roths Satz --- eine Abstandsschranke in $\Ztwo$ ---
    liegt außerhalb der bekannten Theorie.
\end{enumerate}

\medskip
Dieses Dokument benennt diese Lücke als das \textbf{Kernproblem der
arithmetischen Dynamik für Collatz}. Das neue Forschungsfeld ist jetzt
klar abgegrenzt.
\end{tcolorbox}

\subsection{Referenzen}

\begin{itemize}
  \item T.\ Tao: \textit{Almost all orbits of the Collatz map attain almost
    bounded values}. Forum Math.\ Pi \textbf{10} (2022), e12.
  \item J.\ F.\ C.\ Kingman: \textit{The ergodic theory of subadditive stochastic
    processes}. J.\ Roy.\ Stat.\ Soc.\ B \textbf{30} (1968), 499--510.
  \item K.\ F.\ Roth: \textit{Rational approximations to algebraic numbers}.
    Mathematika \textbf{2} (1955), 1--20.
  \item M.\ Li, P.\ Vitányi: \textit{An Introduction to Kolmogorov Complexity
    and Its Applications}. Springer, 4.\ Aufl., 2019.
  \item Ergänzende Dokumente: \textit{collatz\_uniformitaet.tex},
    \textit{collatz\_lte\_dichte.tex}, \textit{collatz\_kingman.tex},
    \textit{collatz\_schlussartikel.tex}.
\end{itemize}

\end{document}
